\begin{abstract}
Language models let attackers change the wording, generator or language of a phishing campaign at almost no cost, while the verified labels a gateway detector needs to adapt arrive late and take analyst time. We study when a fixed label budget should be spent. SLIDE detects shifts of the unlabelled message distribution with Holm-corrected Kolmogorov--Smirnov tests on random projections of a frozen transformer representation and on the detector score, gated by a sliced-Wasserstein effect size. A harmonic level schedule bounds the probability of any false alarm per reference epoch, and each alarm buys the labels of the last 500 messages for any learner. On a 36{,}000-message public stream with human, LLM-generated, paraphrased, PhishFuzzer and machine-translated Spanish segments (ten realisations), SLIDE reached MCC {{main|SLIDE|mcc:.3f}} with {{main|SLIDE|labels:,.0f}} labels. The same number of labels gave {{main|Budget-Uniform-MLP|mcc:.3f}} when spread uniformly, {{main|Budget-Uncertainty-MLP|mcc:.3f}} with uncertainty sampling and {{main|Budget-Random|mcc:.3f}} when bought at random times, against {{main|Budget-Oracle|mcc:.3f}} for an oracle that knows the boundaries; the ordering held for a TF-IDF learner and under gradual, variable-length and permuted streams. An online TF-IDF learner given all {{main|Online-TFIDF-SGD|labels:,.0f}} labels stayed more accurate ({{main|Online-TFIDF-SGD|mcc:.3f}}) unless labels were delayed by eight batches or more. No false alarm occurred on 20 stationary streams. The limits are measured as well: a posterior change with unchanged inputs needed 25 audit labels per batch to be seen, and a benign prefix defeated the 128-token encoder unless messages were scored in windows, which raised the false-positive rate.
\end{abstract}

\begin{keyword}
Phishing detection \sep Concept drift \sep Label budget \sep Sliced Wasserstein distance \sep Large language models \sep Data streams
\end{keyword}
